\scanpage{76}
\mspage{12}
% archive foliation: 78
Au niveau des multiplicités modulaires \inserted{(analytiques réelles.)}
on trouve sûrement un morphisme \inserted{\unsure{pour celles-ci|entre celles-ci}}
\deleted{$\unsure{MB_{g,\nu}}$}
\diagram{\begin{tikzcd}[ampersand replacement=\&]
MDB_{g,\nu} \arrow[r] \& M_{\mathrm{st}}N\deleted{\unsure{{}_{g,\nu}}}D\deleted{\unsure{{}_{g}}}{}_{g,\nu}
\end{tikzcd}}
où $MDB_{g,\nu}$ désigne la multiplicité modulaire
(M) pour les surfaces holomorphes compactes
``découpées'' (D) à bord (B) de type $g,\nu$, et
$M_{\mathrm{st}}ND_{g,\nu}$, celle pour les courbes algébriques
stables (de type loc. constant) \inserted{de type $g,\nu$} à structure de
propagation \unsure{unnodale|nodale} (N) et à système \inserted{adm.} de
disques (D) autour des pts marqués.

Cette flèche est compatible avec les
flèches canoniques vers la cat. des
types combinatoires\inserted{-toriques} de tels objets
(diagrammes MD (Mumford-Deligne en
``\unsure{modulaire|modulaires}'') (\deleted{\unsure{avec poids}}\inserted{incluant} une pondération
des sommets) à structures toriques
associées aux arêtes \ldots\ donnant lieu à un
diagramme commutatif
